\section{Assembly and same-space consequences}\label{sec:assembly}

The construction now supplies both inputs needed after the absorption
criterion: a lifted correction and the absence of linear \(c_0\).

\begin{proof}[Proof of Theorem~\ref{thm:main}]
Lemma~\ref{lem:completion} gives a separable real Banach space \(Z\).
Proposition~\ref{prop:no-c0} excludes a linearly isomorphic copy of
ordinary \(c_0\) in \(Z\). Propositions~\ref{prop:global} and
\ref{prop:onto} give a bi-Lipschitz bijection
\[
                      g:E\oplus_\infty U\longrightarrow U
\]
with \(E=c_0(\N)\) and lower constant
\(c_g=1/5-24\eta>0\). Lemmas~\ref{lem:Q} and \ref{lem:K} give a
bounded linear \(Q:Z\to U\) and a Lipschitz \(K:E\oplus_\infty U\to Z\)
with \(\Lip(K)\leq3\) and \(QK=g-P\). Applying
Lemma~\ref{lem:absorption} yields the onto bi-Lipschitz map
\[
             F:Z\oplus_\infty E\to Z,\qquad
             F(b,t)=b+K(t,Qb),
\]
with the explicit inverse \eqref{eq:absorb-inverse}. Its inverse is the
bi-Lipschitz bijection in the opposite direction
\(Z\to Z\oplus_\infty c_0\). This proves the theorem.
\end{proof}

\begin{samepage}
\begin{corollary}\label{cor:same-space}
For this same separable real Banach space \(Z\):
\begin{enumerate}
\item \(Z\oplus_\infty c_0\) and \(Z\) are bi-Lipschitz equivalent
but not linearly isomorphic.
\item \(c_0\) embeds bi-Lipschitzly into \(Z\), although \(Z\)
contains no linear copy of \(c_0\).
\item Every separable metric space embeds bi-Lipschitzly into \(Z\).
\end{enumerate}
\end{corollary}
\end{samepage}

\begin{proof}
The direct summand \(\{0\}\oplus c_0\) is a linear isometric copy of
\(c_0\). A bounded linear isomorphism from the direct sum onto \(Z\)
would carry it to a closed linear copy in \(Z\), contradicting
Proposition~\ref{prop:no-c0}. The restriction \(t\mapsto F(0,t)\)
is a bi-Lipschitz embedding, proving the second statement.
Aharoni's theorem \cite{Aharoni74} embeds every separable metric space
bi-Lipschitzly into \(c_0\); composition with this restriction gives
the third statement.
\end{proof}
